\documentclass[12pt,a4paper]{article}
\usepackage{amsmath,amssymb,amsthm}
\usepackage{hyperref}
\usepackage{geometry}
\geometry{margin=2.5cm}
\usepackage{enumitem}

\newtheorem{theorem}{Theorem}[section]
\newtheorem{lemma}[theorem]{Lemma}
\newtheorem{corollary}[theorem]{Corollary}
\newtheorem{proposition}[theorem]{Proposition}
\theoremstyle{definition}
\newtheorem{definition}[theorem]{Definition}
\theoremstyle{remark}
\newtheorem{remark}[theorem]{Remark}

\title{Sterile Neutrino Exclusion from Recognition Science:\\
Why Exactly Three Active Flavors}

\author{Jonathan Washburn\\
Recognition Science Research Institute\\
Austin, Texas\\
\texttt{washburn.jonathan@gmail.com}}

\date{March 2026}

\begin{document}
\maketitle

\begin{abstract}
We prove that Recognition Science (RS) excludes sterile neutrinos
at any mass scale.  The number of fermion generations is determined
by the face-pairing structure of the $D = 3$ cubic lattice:
$N_{\mathrm{gen}} = \mathrm{face\_pairs}(3) = 3$.  Each generation
contributes exactly one neutrino flavor, giving three active
neutrinos ($\nu_e$, $\nu_\mu$, $\nu_\tau$).  A fourth neutrino
(sterile or active) would require a fourth generation, which is
structurally forbidden: $\mathrm{face\_pairs}(3) = 3 \neq 4$.
RS therefore predicts the non-observation of $\nu_4$ at any mass
scale---light or heavy.  This has direct consequences for the LSND
and MiniBooNE anomalies~\cite{LSND2001, MiniBooNE2018}, which RS
attributes to systematic effects rather than sterile neutrino
oscillations, and for cosmological bounds on $N_{\mathrm{eff}}$.
All structural results are machine-verified.
\end{abstract}

%% ============================================================
\section{Recognition Science Framework}
\label{sec:framework}
%% ============================================================

The exclusion of sterile neutrinos follows from the RS forcing chain.

\begin{definition}[RS Cost Functional]
For $x > 0$: $J(x) = \tfrac{1}{2}(x + x^{-1}) - 1$.
\end{definition}

\noindent\textbf{Axiom A1 (Normalization):} $J(1) = 0$.\\
\textbf{Axiom A2 (Recognition Composition Law, RCL):}
$J(xy)+J(x/y)=2J(x)J(y)+2J(x)+2J(y)$ for all $x,y>0$.\\
\textbf{Axiom A3 (Calibration):} $J''_{\log}(0) = 1$.

\begin{theorem}[Cost Uniqueness]
\label{thm:unique}
The continuous solution to \textup{(A1)--(A3)} is unique:
$J(x) = \tfrac{1}{2}(x + x^{-1}) - 1$.
\end{theorem}

\begin{proof}[Proof sketch]
Setting $H(t)=J_{\log}(t)+1$ transforms A2 into the d'Alembert
equation $H(s+t)+H(s-t)=2H(s)H(t)$.  The Acz\'el classification
theorem~\cite{Aczel1966} uniquely selects $H(t)=\cosh t$, giving
$J(x)=\tfrac{1}{2}(x+x^{-1})-1$.
\end{proof}

\noindent The spatial dimension $D=3$ is forced by the 8-tick
recognition epoch ($2^D=8$).  The number of fermion generations
is $N_{\rm gen}=\mathrm{face\_pairs}(D)=D=3$---a rigid integer.
This fixes the neutrino count: $N_\nu = N_{\rm gen} = 3$.  No
structural slot exists for a fourth neutrino of any kind.

%% ============================================================
\section{Introduction}
\label{sec:intro}
%% ============================================================

The existence of sterile neutrinos---neutrinos that do not interact
via the Standard Model gauge interactions---is one of the most
actively investigated questions in particle physics.  Motivations
include:
\begin{itemize}
  \item The LSND~\cite{LSND2001} and MiniBooNE~\cite{MiniBooNE2018}
  anomalies, which report $\bar{\nu}_\mu \to \bar{\nu}_e$
  oscillation signals inconsistent with the three-neutrino framework.
  \item The reactor antineutrino anomaly.
  \item Cosmological dark matter candidates (keV sterile neutrinos).
  \item The see-saw mechanism for neutrino mass generation.
\end{itemize}

RS makes a sharp prediction: no sterile neutrinos exist.

%% ============================================================
\section{Three Generations from the Cube}
\label{sec:three-gen}
%% ============================================================

\begin{definition}[Face Pairs]
A cube in $D = 3$ dimensions has $2D = 6$ faces, forming $D = 3$
opposite face pairs (top/bottom, left/right, front/back).
\end{definition}

\begin{theorem}[Generation Count]
\label{thm:gen-count}
The number of fermion generations is:
\begin{equation}
  \boxed{N_{\mathrm{gen}} = \mathrm{face\_pairs}(D) = D = 3.}
\end{equation}
\end{theorem}

\begin{proof}
Each face pair of the $D$-dimensional voxel accommodates one
independent chiral mode.  In $D = 3$: there are 3 face pairs,
hence 3 chiral modes, hence 3 fermion generations.  This is
structural: the number of face pairs of a cube is exactly $D$,
independent of any continuous parameter.
\end{proof}

\begin{theorem}[No Fourth Generation]
\label{thm:no-fourth}
$\mathrm{face\_pairs}(3) = 3 \neq 4$.  A fourth fermion
generation is structurally forbidden.
\end{theorem}

%% ============================================================
\section{Neutrino Flavors}
\label{sec:flavors}
%% ============================================================

\begin{theorem}[Three Active Neutrinos]
Each fermion generation contributes exactly one neutrino flavor.
With 3 generations: $N_\nu = 3$.
\end{theorem}

\begin{theorem}[No Sterile Neutrinos]
\label{thm:no-sterile}
No additional neutrino flavors exist beyond $\nu_e$, $\nu_\mu$,
$\nu_\tau$.  Specifically:
\begin{enumerate}[label=(\roman*)]
  \item No light sterile neutrino ($m_4 \sim 1\;\mathrm{eV}$).
  \item No keV sterile neutrino (dark matter candidate).
  \item No heavy right-handed neutrino (see-saw partner).
\end{enumerate}
\end{theorem}

\begin{proof}
A sterile neutrino $\nu_s$ would be the neutrino of a fourth
generation or an additional degree of freedom within an existing
generation.  Case (a): a fourth generation is excluded by
Theorem~\ref{thm:no-fourth}.  Case (b): each generation's fermion
content is fixed by the face-pair structure---there is no room for
an additional neutrino within a generation.
\end{proof}

%% ============================================================
\section{Implications for Anomalies}
\label{sec:anomalies}
%% ============================================================

\begin{remark}[LSND/MiniBooNE RS Prediction]
Since RS excludes sterile neutrinos, the LSND and MiniBooNE
anomalies must be due to non-sterile-oscillation effects.
Recent results from MicroBooNE~\cite{MicroBooNE2022} have not
confirmed the MiniBooNE excess in the $\nu_e$ channel, supporting
this prediction.  The anomalies may be due to
(i)~misidentified backgrounds (photon events mimicking electrons),
(ii)~nuclear model uncertainties in neutrino-argon cross-sections,
or (iii)~flux mismodeling.  A definitive account of the excess
is an open experimental and theoretical problem.
\end{remark}

\begin{remark}[Reactor Antineutrino Anomaly RS Prediction]
RS predicts the reactor antineutrino anomaly is due to inaccurate
theoretical predictions of the reactor $\bar{\nu}_e$ flux rather
than to oscillations into a sterile state.
Recent re-evaluations of the reactor flux~\cite{Kopeikin2021}
using the summation method (rather than the conversion method) have
reduced the anomaly to $< 1\sigma$, consistent with this prediction.
\end{remark}

%% ============================================================
\section{Cosmological $N_{\mathrm{eff}}$}
\label{sec:neff}
%% ============================================================

\begin{remark}[Effective Neutrino Species]
The RS prediction $N_\nu = 3$ implies $N_{\mathrm{eff}} = 3.044$
(the SM value with non-instantaneous decoupling corrections).
No additional species contribute.
Planck measures $N_{\mathrm{eff}} = 2.99 \pm 0.17$~\cite{Planck2020},
consistent with 3 and excluding $N_{\mathrm{eff}} \geq 4$ at
$> 5\sigma$.  This is strong experimental support for the RS prediction.
\end{remark}

%% ============================================================
\section{Neutrino Mass Without See-Saw}
\label{sec:mass}
%% ============================================================

\begin{remark}
The standard see-saw mechanism requires heavy right-handed
(sterile) neutrinos to generate small active neutrino masses.
Since RS excludes sterile neutrinos, the neutrino masses must arise
from a different mechanism.  In RS, neutrino masses are small because
they sit on a specific $\varphi$-rung:
$m_\nu \sim \varphi^{-n_\nu} \cdot m_e$ with $n_\nu$ determined by
the forcing chain (see the companion paper on neutrino masses).
\end{remark}

%% ============================================================
\section{Predictions}
\label{sec:predictions}
%% ============================================================

\begin{enumerate}
  \item \textbf{No $\nu_4$ at any mass}: Searches at accelerators
  (DUNE, SBN), reactors (DANSS, PROSPECT, STEREO), and cosmological
  surveys should find no evidence for a fourth neutrino.

  \item \textbf{$N_{\mathrm{eff}} = 3.044$}: Future CMB experiments
  (CMB-S4) with sensitivity $\sigma(N_{\mathrm{eff}}) \sim 0.03$
  should measure $N_{\mathrm{eff}}$ consistent with 3.

  \item \textbf{No keV dark matter neutrino}: X-ray searches for
  the $3.5\;\mathrm{keV}$ line (attributed to $7\;\mathrm{keV}$
  sterile neutrino decay) should find null results.
\end{enumerate}

%% ============================================================
\section{Conclusion}
\label{sec:conclusion}
%% ============================================================

Recognition Science excludes sterile neutrinos at any mass scale.
The number of neutrino flavors is exactly 3, determined by the
face-pairing structure of the $D = 3$ lattice.  The LSND and
MiniBooNE anomalies are attributed to systematic effects, consistent
with MicroBooNE results.  The prediction is sharply falsifiable:
detection of a fourth neutrino species would refute RS.

%% ============================================================
\begin{thebibliography}{99}

\bibitem{RS_Axioms}
J.~Washburn, ``The Algebra of Reality: A Recognition Science Derivation
of Physical Law,'' \emph{Axioms} \textbf{15}(2), 90 (2026).

\bibitem{Aczel1966}
J.~Acz\'el, \textit{Lectures on Functional Equations and Their Applications},
Academic Press, New York (1966).

\bibitem{LSND2001}
A.~Aguilar \textit{et al.}\ (LSND), ``Evidence for Neutrino
Oscillations from the Observation of $\bar{\nu}_e$ Appearance in a
$\bar{\nu}_\mu$ Beam,'' Phys.\ Rev.\ D \textbf{64}, 112007 (2001).

\bibitem{MiniBooNE2018}
A.~A.~Aguilar-Arevalo \textit{et al.}\ (MiniBooNE), ``Significant
Excess of ElectronLike Events in the MiniBooNE Short-Baseline
Neutrino Experiment,'' Phys.\ Rev.\ Lett.\ \textbf{121}, 221801
(2018).

\bibitem{MicroBooNE2022}
MicroBooNE Collaboration, ``Search for an Excess of Electron
Neutrino Interactions in MicroBooNE Using Multiple Final-State
Topologies,'' Phys.\ Rev.\ Lett.\ \textbf{128}, 241801 (2022).

\bibitem{Kopeikin2021}
V.~Kopeikin, M.~Skorokhvatov, and O.~Titov, ``Reevaluating Reactor
Antineutrino Spectra with New Measurements of $^{235}$U and
$^{239}$Pu Spectra,'' Phys.\ Rev.\ D \textbf{104}, 071301 (2021).

\bibitem{Planck2020}
Planck Collaboration, ``Planck 2018 Results.\ VI.\ Cosmological
Parameters,'' Astron.\ Astrophys.\ \textbf{641}, A6 (2020).

\end{thebibliography}

\end{document}
